% Image credits for the photographs and AI illustrations of Book 2 (High School).
% Machine-readable license records live in images/book2/CREDITS.md; keep the
% two files in sync. The Book 2 writer owns this file: add one \item per
% photograph (inside the itemize, by chapter) as photographs land.
% Spanish edition: every author, licence and source of the English list is
% kept exactly; only the connecting words and the chapter titles are translated.
\chapter*{Créditos de las imágenes}

Los dibujos, mapas y gráficos de este libro son originales; los mapas y
gráficos trazados a partir de datos publicados citan su fuente en el pie de
figura o en esta lista. Las fotografías se utilizan con sus respectivas
licencias libres, con nuestro agradecimiento a sus autores; los enlaces
completos a las fuentes y las URL de las licencias de cada fotografía
figuran en \texttt{images/book2/CREDITS.md}, en el repositorio del libro.

\bigskip
Junto a las fotografías y a los dibujos originales, algunas figuras son
ilustraciones generadas por IA, producidas con la generación de imágenes de
OpenAI a partir de instrucciones escritas por los editores del libro, y
revisadas en cuanto a su exactitud científica antes de incluirlas. Solo
muestran escenas genéricas, nunca un lugar real, una persona real ni un
ejemplar real. La instrucción completa de cada imagen generada está
registrada en \texttt{images/book2/ai/PROMPTS.md}, en el repositorio.

\bigskip
\noindent Fotografías y datos, por capítulo:
\begin{itemize}\small
  \item Cap.~1 (medir la Tierra): la Tierra vista desde el Apolo~17; fotografía de la NASA, dominio público (Wikimedia Commons).
  \item Cap.~1: un círculo repetidor, lámina de J.-B. Delambre, \emph{Base du système métrique décimal}, vol.~2 (1807); dominio público (Wikimedia Commons, a partir de BnF Gallica).
  \item Cap.~1: longitudes de grado calculadas sobre el elipsoide WGS~84 (conjunto de datos EPSG).
  \item Cap.~2 (un planeta habitable): Venus desde la Mariner~10 (PIA23791), NASA/JPL-Caltech, dominio público; el monte Sharp desde el róver Curiosity (PIA16769), NASA/JPL-Caltech/MSSS, dominio público (Wikimedia Commons).
  \item Cap.~2: datos planetarios de las fichas técnicas de la NASA; límites de fase del agua de las publicaciones de la IAPWS SR1-86 y R14-08; límites de la zona habitable de Kopparapu y otros (2013); Próxima~b del NASA Exoplanet Archive.
  \item Cap.~3 (el balance energético de la Tierra): charcas de fusión sobre el hielo marino ártico (misión ICESCAPE, 2011); fotografía del NASA Goddard Space Flight Center, dominio público (Wikimedia Commons).
  \item Cap.~3: valores del balance energético del Sexto Informe de Evaluación del IPCC, Grupo de Trabajo~I, figura~7.2; albedos de la University of Colorado Boulder (MOSAiC) y del National Snow and Ice Data Center; transporte de calor de Trenberth y Caron (2001); producción primaria neta de Field y otros (1998); datos termoquímicos del NIST Chemistry WebBook.
  \item Cap.~4 (por qué se forman las nubes): cúmulos de buen tiempo, fotografía de Ralph F. Kresge, NOAA; el yunque de un cumulonimbo, NOAA National Severe Storms Laboratory; ambas de dominio público (Wikimedia Commons).
  \item Cap.~4: saturación calculada a partir de las publicaciones de la IAPWS; tamaños de gotas del NASA Jet Propulsion Laboratory.
  \item Cap.~5 (tormentas y ciclones tropicales): el huracán Isabel desde la Estación Espacial Internacional, fotografía de Mike Trenchard, NASA Johnson Space Center; un tornado en Kansas, fotografía de Sean Waugh, NOAA National Severe Storms Laboratory; rayos, fotografía de C.~Clark, NOAA National Severe Storms Laboratory; todas de dominio público (Wikimedia Commons).
  \item Cap.~5: trayectorias de ciclones de IBTrACS versión~4r01 (Gahtan y otros, 2024; Knapp y otros, 2010; NOAA National Centers for Environmental Information); categorías de huracanes del NOAA National Hurricane Center; energía de los huracanes y condiciones de formación del NOAA Atlantic Oceanographic and Meteorological Laboratory; frecuencia de rayos de Christian y otros (2003); datos sobre tornados del NOAA Storm Prediction Center.
  \item Cap.~6 (el océano estratificado): una roseta de muestreo, fotografía de Pablo Clemente-Colon, NOAA Photo Library; una boya perfiladora lanzada desde un buque oceanográfico, NOAA Global Ocean Monitoring and Observing Program; clorofila oceánica 2002--2010, Jesse Allen y Robert Simmon, NASA Earth Observatory; todas de dominio público (Wikimedia Commons).
  \item Cap.~6: perfiles oceánicos del World Ocean Atlas 2023 (Reagan y otros, 2024; NOAA National Centers for Environmental Information); densidad del agua de mar de la ecuación de estado internacional (Fofonoff y Millard, 1983); volumen del océano de Charette y Smith (2010); datos sobre las boyas de la Argo Program Office; afloramiento del NOAA National Ocean Service y de Kessler (2006).
  \item Cap.~7 (el suelo, un recurso de formación lenta): un perfil de laterita sobre gneis en Kerala, fotografía de Werner Schellmann, CC~BY-SA~2.5; saprolita en el centro de Madagascar, fotografía de Diorit, CC~BY-SA~3.0; un perfil de suelo, USDA Natural Resources Conservation Service, dominio público (todas Wikimedia Commons).
  \item Cap.~7: ritmos de producción y de erosión del suelo de Montgomery (2007); fórmulas de minerales de la Mineralogy Database (webmineral.com); espesor de la saprolita del USGS Ground Water Atlas of the United States.
  \item Cap.~8 (cuencas de drenaje, ríos e inundaciones): un hidrólogo midiendo un caudal y una estación de aforo, U.S. Geological Survey, dominio público (Wikimedia Commons).
  \item Cap.~8: máximos anuales y caudal diario del río Potomac cerca de Washington, U.S. Geological Survey National Water Information System, estación 01646500; caudal y cuenca del Misisipi de Kammerer (1990, USGS); caudal del Amazonas de Callède y otros (2010); seguridad ante el agua de las inundaciones del NOAA National Weather Service.
  \item Cap.~9 (los recursos de agua dulce): el mar de Aral en 2000 y 2018, NASA Earth Observatory (MODIS/Terra); campos de regadío cerca de Sublette, Kansas, NASA (ASTER/Terra), dominio público (Wikimedia Commons).
  \item Cap.~9: reservas mundiales de agua según Shiklomanov (1993), U.S. Geological Survey Water Science School; el ciclo global del agua según Trenberth y otros (2007), tal como se publica en Trenberth y Asrar (2014); edades de las aguas subterráneas de Nubia de Sturchio y otros (2004); cambios del acuífero de las High Plains de McGuire (2017, USGS); balance del mar de Aral de Gaybullaev y otros (2012) y de Micklin (2007).
  \item Cap.~10 (de la deriva continental a la expansión del fondo oceánico): el buque de perforación \emph{JOIDES Resolution}, Ocean Drilling Program, Texas A\&M University, dominio público (Wikimedia Commons).
  \item Cap.~10: la escala de tiempo de polaridad de Cande y Kent (1995); los resultados de perforación de la Leg~3 del Deep Sea Drilling Project (1970); la edad del fondo oceánico de Seton y otros (2020), EarthByte, licencia CC~BY~4.0.
  \item Cap.~11 (las ondas sísmicas revelan el interior): instalación de un sismómetro temporal en Puerto Rico; fotografía de Thomas Pratt, U.S. Geological Survey, dominio público (Wikimedia Commons).
  \item Cap.~11: el Preliminary Reference Earth Model de Dziewonski y Anderson (1981); tiempos de recorrido empíricos de las ondas P y S del Australian Passive Seismic Server (AuSPASS).
  \item Cap.~12 (litosfera y astenosfera): un iceberg, NOAA, dominio público (USGS Water Science School).
  \item Cap.~12: flujo de calor de Davies y Davies (2010), licencia CC~BY~3.0; provincias de flujo de calor de Roy, Blackwell y Decker (1972); el solidus del manto de Hirschmann (2000); espesor y densidad de la corteza de Christensen y Mooney (1995); altitudes de ETOPO1 (NOAA).
  \item Cap.~13 (las placas y sus movimientos): una estación de posicionamiento por satélite en la cordillera de las Cascadas; fotografía de Dave Sherrod, U.S. Geological Survey, dominio público (Wikimedia Commons). Relieve sombreado del fondo del Pacífico Norte, National Geophysical Data Center, NOAA, dominio público (Wikimedia Commons).
  \item Cap.~13: polos de rotación de las placas NNR-MORVEL56 de Argus, Gordon y DeMets (2011); posiciones de estaciones del Nevada Geodetic Laboratory (Blewitt, Hammond y Kreemer, 2018); edades de volcanes de Clague y Dalrymple (1987), U.S. Geological Survey; límites de placas de Bird (2003); tierras emergidas de Natural Earth.
  \item Cap.~14 (las dorsales oceánicas): lava almohadillada en el rift de las Galápagos, NOAA Okeanos Explorer Program, dominio público; un fumador negro en la dorsal de Juan de Fuca, University of Washington y NOAA, dominio público; gabro, fotografía de Darla Sondrol (GeoDIL), CC0; peridotita en basalto, fotografía de James St.~John, CC~BY~2.0 (todas Wikimedia Commons).
  \item Cap.~14: profundidades del fondo oceánico de ETOPO1 (NOAA) y edades de Seton y otros (2020), EarthByte, licencia CC~BY~4.0; leyes de profundidad de Parsons y Sclater (1977); el solidus del manto de Hirschmann (2000); polos de rotación NNR-MORVEL56 de Argus, Gordon y DeMets (2011).
  \item Cap.~15 (las zonas de subducción): esquisto azul y eclogita, fotografías de James St.~John, CC~BY~2.0; daños del tsunami en Kesennuma, Japón, 2011, fotografía de Dylan McCord, U.S. Navy, dominio público (todas Wikimedia Commons).
  \item Cap.~15: terremotos del catálogo de terremotos del USGS (ComCat); límites de placas y velocidades de convergencia de Bird (2003); edades del fondo oceánico de Seton y otros (2020), EarthByte, licencia CC~BY~4.0; velocidades de placas NNR-MORVEL56 de Argus, Gordon y DeMets (2011); tierras emergidas de Natural Earth.
  \item Cap.~16 (los minerales): cuarzo, ortosa, biotita y olivino, fotografías de Shannon Heinle, Darla Sondrol, Dexter Perkins y Nessa Eull (Geology Digital Image Library), CC0; basalto en lámina delgada, fotografías de Chmee2, CC~BY-SA~3.0; cuarzo de un granito en lámina delgada, fotografía de Michael C.~Rygel, CC~BY-SA~3.0 (todas Wikimedia Commons).
  \item Cap.~16: la composición de la corteza continental de Rudnick y Gao (2014); pesos atómicos estándar de IUPAC/CIAAW; fórmulas y densidades de minerales de webmineral.com; densidades de rocas graníticas de Sierra Nevada del U.S. Geological Survey (Open-File Report 91-570-A).
  \item Cap.~17 (las rocas ígneas): columnas de basalto de la Calzada del Gigante, fotografía de Adam Bishop, CC~BY-SA~3.0; El Capitan, Yosemite, U.S. Geological Survey, CC0; basalto, andesita, riolita, diorita, granito, obsidiana, una roca porfídica y dos láminas delgadas, fotografías de Nessa Eull, Twyla Baker, Joseph H.~Hartman, Mike Davis, Dexter Perkins y Darla Sondrol (Geology Digital Image Library), CC0 (todas Wikimedia Commons).
  \item Cap.~17: enfriamiento de diques y filones capa según Delaney (1986), U.S. Geological Survey Open-File Report 86-444; composiciones de magmas y orden de cristalización de los apuntes de geología física de S.~A.~Nelson y de S.~Earle, \emph{Physical Geology} (CC~BY~4.0); la biografía de Bowen de la historia del Geophysical Laboratory de Carnegie Science.
  \item Cap.~18 (carbón, petróleo y gas): plantas fósiles en lutita rica en carbono, lignito, carbón bituminoso y antracita, fotografías de James St.~John, CC~BY~2.0; una capa de lignito en Rumanía, U.S. Geological Survey, dominio público (todas Wikimedia Commons).
  \item Cap.~18: ventana del petróleo, formación del carbón y trampas de S.~Earle, \emph{Physical Geology} (CC~BY~4.0); contenido en carbono del carbón de la U.S. Energy Information Administration; el sistema petrolífero según Magoon y Schmoker (U.S. Geological Survey, 2000); las cuencas de rift y salinas del oeste de África Central de la U.S. Geological Survey Fact Sheet 2023-3028.
  \item Cap.~19 (la corteza continental): el gneis de Acasta, fotografía de Pedroalexandrade, CC~BY-SA~3.0; circones de Jack Hills, imágenes de A.~Cavosie y J.~Valley, dominio público; gneises y granitos del escudo canadiense, fotografía de Chris Light, CC~BY-SA~4.0 (todas Wikimedia Commons).
  \item Cap.~19: composición de la corteza de Rudnick y Gao (2014); espesor de la corteza de Christensen y Mooney (1995), de White, McKenzie y O'Nions (1992) y de Karplus y otros (2019); extensión de los continentes de Cogley (1984); edades del fondo oceánico de Seton y otros (2020), EarthByte, licencia CC~BY~4.0.
  \item Cap.~20 (la formación de las montañas): el monte Everest desde el transbordador espacial Columbia, NASA, dominio público; capas de roca plegadas en la garganta del Kali Gandaki, fotografía de Jean-Marie Hullot, CC~BY~2.0; el cabalgamiento de Glarus, fotografía de Christian Heine, CC~BY-SA~3.0; lava almohadillada de la ofiolita de Semail, fotografía de Lionel Duchoiselle, CC~BY-SA~3.0 (todas Wikimedia Commons).
  \item Cap.~20: posiciones pasadas de la India y de Lhasa de los modelos de placas de Müller y otros (2019, 2022), calculadas con el GPlates Web Service (EarthByte); relieves altos de ETOPO1; terremotos del catálogo del USGS; límites de placas de Bird (2003); tierras emergidas de Natural Earth.
  \item Cap.~21 (las rocas metamórficas): esquistosidad pizarrosa al microscopio, fotografía de Geoloshe, CC0; micaesquisto con granate de Siros, fotografía de Graeme Churchard, CC~BY~2.0; esquisto con distena y gneis de cuarzo y biotita, fotografías de James St.~John, CC~BY~2.0 (todas Wikimedia Commons).
  \item Cap.~21: los campos de estabilidad de la distena, la andalucita y la sillimanita de Holdaway y Mukhopadhyay (1993).
  \item Cap.~22 (el magmatismo de las zonas de subducción): la columna eruptiva del monte Santa Helena, 18 de mayo de 1980, fotografía de Donald A. Swanson, USGS, dominio público; granodiorita, andesita y Sierra Nevada, fotografías de James St.~John, CC~BY~2.0 (todas Wikimedia Commons).
  \item Cap.~22: terremotos del catálogo del USGS; volcanes del Smithsonian Global Volcanism Program; límites de placas de Bird (2003); tierras emergidas de Natural Earth; el solidus de la peridotita seca de Hirschmann (2000), con las densidades del PREM.
  \item Cap.~23 (la destrucción de las montañas): los montes Apalaches desde la Estación Espacial Internacional, fotografía de astronauta ISS033-E-22378, NASA, dominio público; conglomerado de cantos de cuarzo, fotografía de James St.~John, CC~BY~2.0 (ambas Wikimedia Commons).
  \item Cap.~23: ritmos de denudación de Portenga y Bierman (2011); cargas de sedimentos del Misisipi de Meade y Moody (2010); la constante de decrecimiento de Ahnert (1970); las Varden Hills, la región inventada de la serie.
  \item Cap.~24 (el calor interno de la Tierra): la central geotérmica de Nesjavellir, Islandia, fotografía de Gretar Ívarsson, editada por Fir0002, dominio público (Wikimedia Commons).
  \item Cap.~24: flujo de calor de Davies y Davies (2010); calor radiogénico de la medida de geoneutrinos de KamLAND (2011) y de la liberación actual de calor de los cuatro isótopos; producción de calor de las rocas de Vilà y Fernández (2010).
  \item Cap.~25 (la datación radiométrica y la edad de la Tierra): William Thomson, lord Kelvin, 1906, fotografía de Russell and Sons, Londres, Smithsonian Institution Libraries, sin restricciones de derechos de autor conocidas; el meteorito metálico de Canyon Diablo, fotografía de Geoffrey Notkin, CC~BY-SA~2.5 (ambas Wikimedia Commons).
  \item Cap.~25: isótopos del plomo de los meteoritos de Patterson (1956); la constante de desintegración del rubidio 87 recomendada por la IUGS y la IUPAC (2022); la historia de las estimaciones a partir de Holmes (1913, 1927), Patterson (1956) y el U.S. Geological Survey; las Varden Hills, la región inventada de la serie.
  \item Cap.~26 (la Tierra primitiva y su atmósfera): formación de hierro bandeado, garganta de Dales, fotografía de Graeme Churchard, CC~BY~2.0; estromatolitos en la bahía Shark, fotografía de Paul Harrison, CC~BY-SA~3.0; capas rojas cerca de Sedona, fotografía de Wendy from Pennsylvania, CC~BY-SA~2.0 (todas Wikimedia Commons).
  \item Cap.~26: la luminosidad del Sol joven de Gough (1981); límites y colores de los eones de la Tabla Cronoestratigráfica Internacional.
  \item Cap.~27 (leer los climas del pasado): burbujas de aire en hielo antártico, fotografía de CSIRO Atmospheric Research, CC~BY~3.0; foraminíferos planctónicos, fotografía de Marci Robinson, USGS, dominio público; granos de polen, Dartmouth Electron Microscope Facility, dominio público (todas Wikimedia Commons).
  \item Cap.~27: temperatura de EPICA Domo~C (Jouzel y otros, 2007) y dióxido de carbono de los testigos de hielo antárticos (Bereiter y otros, 2015), vía NOAA Paleoclimatology; el apilamiento LR04 (Lisiecki y Raymo, 2005); relieve ETOPO1 (NOAA); NASA GISS (GISTEMP v4); periodos de referencia del Sexto Informe de Evaluación del IPCC, Grupo de Trabajo~I, capítulo~2.
  \item Cap.~28 (el cambio climático): el observatorio de Mauna Loa; fotografía de Brian Vasel, NOAA, dominio público (Wikimedia Commons).
  \item Cap.~28: Charles David Keeling y John Chin en Mauna Loa, 1997; fotografía de John Miller, NOAA, dominio público (Wikimedia Commons).
  \item Cap.~28: datos del NOAA Global Monitoring Laboratory y de la Scripps Institution of Oceanography (dióxido de carbono), de NASA GISS (GISTEMP v4), del NOAA Laboratory for Satellite Altimetry (nivel del mar), de Bereiter y otros (2015) vía NOAA Paleoclimatology (testigos de hielo) y del Sexto Informe de Evaluación del IPCC, Grupo de Trabajo~I (forzamientos, atribución, proyecciones).
\end{itemize}
